
\begin{table}[H]
   \caption{\label{tab:scheduling} Cesareans dip on weekends and holidays, more so in the private sector}
   \centering
\small\setlength{\tabcolsep}{5pt}
   \begin{tabular}{lcc}
      \toprule
      Dependent Variable: & \multicolumn{2}{c}{rate}\\
                       & Public          & Private \\   
      Model:           & (1)             & (2)\\  
      \midrule
      \emph{Variables}\\
      Weekend          & -0.0659$^{***}$ & -0.0801$^{***}$\\   
                       & (0.0024)        & (0.0043)\\   
\addlinespace[2pt]
      National holiday & -0.0343$^{***}$ & -0.0512$^{***}$\\   
                       & (0.0014)        & (0.0025)\\   
\addlinespace[2pt]
      Eve of rest day  & -0.0043$^{***}$ & -0.0097$^{***}$\\   
                       & (0.0014)        & (0.0017)\\   
\addlinespace[2pt]
      \midrule
      \emph{Fixed-effects}\\
      Municipality     & Yes             & Yes\\  
      Year             & Yes             & Yes\\  
      \midrule
      \emph{Fit statistics}\\
      Observations     & 2,416,556       & 926,770\\  
      R$^2$            & 0.363           & 0.440\\  
\bottomrule
   \end{tabular}
   
   \par \raggedright 
   \footnotesize\textit{Notes:} Municipality-date cells, weighted by births. The dependent variable is the cesarean share of births. \emph{Eve of rest day} is a Friday or the day before a national holiday. Movable holidays (Good Friday, Carnival, Corpus Christi) are included. Standard errors clustered by municipality. \footnotesize Significance: *** p$<$0.01, ** p$<$0.05, * p$<$0.10.
\end{table}


